% ch4.tex -- Module 8 ESSENTIALS edition, Chapter 4: Properties of the Nucleus (IQ4).
% book): same notation, same facts and numbers. ch4_becurve.pdf is copied from
% src/mod8/figures/; the other diagrams are TikZ copied (some shrunk) from the full
% chapter. Dot points M8.IQ4.01 - M8.IQ4.06.
\chapter[Properties of the Nucleus]{Properties of the Nucleus: Radioactivity, Fission and Fusion}\label{ch4:chap}
\covers{M8.IQ4.01, M8.IQ4.02, M8.IQ4.03, M8.IQ4.04, M8.IQ4.05, M8.IQ4.06}

% Chapter macros (prefix chFour, as in the full book):
%   \chFourNuc{A}{Z}{X}  nuclide symbol, A above Z, both right-aligned before X
\newcommand\chFourNuc[3]{\ensuremath{\mathord{%
  \setbox0=\hbox{$\scriptstyle#1$}\setbox2=\hbox{$\scriptstyle#2$}%
  \ifdim\wd0<\wd2 \setbox0=\hbox to\wd2{\hss\box0}\else\setbox2=\hbox to\wd0{\hss\box2}\fi
  {}^{\box0}_{\box2}\mathrm{#3}}}}
\newcommand\chFourAlpha{\chFourNuc{4}{2}{He}}
\newcommand\chFourBeta{\chFourNuc{0}{-1}{e}}
\newcommand\chFourPos{\chFourNuc{0}{+1}{e}}
\newcommand\chFourNeut{\chFourNuc{1}{0}{n}}
\newcommand\chFourGam{\ensuremath{{}^{0}_{0}\gamma}}
\newcommand\chFourAnu{\ensuremath{\overline{\nu}_\mathrm{e}}}
\newcommand\chFourNu{\ensuremath{\nu_\mathrm{e}}}

\begin{summarybox}{In this chapter}
\begin{itemize}
  \item Unstable nuclei: alpha, beta and gamma decay and their radiation.
  \item Half-life, the decay law and activity.
  \item Mass defect, binding energy and the energy released in nuclear reactions.
  \item Fission, chain reactions, reactors and fusion.
\end{itemize}
\end{summarybox}

%% ==================================================================
\section{Radioactive decay: alpha, beta and gamma}\label{ch4:sec:decay}
\covers{M8.IQ4.01, M8.IQ4.04}

\begin{bigidea}
Some nuclei have the wrong mix of protons and neutrons, or are too big. They change by
themselves into more stable nuclei and give out alpha, beta or gamma radiation.
\end{bigidea}

\begin{keywords}
  \kw{Nucleon}{a proton or a neutron (a particle in the nucleus).}
  \kw{Strong nuclear force}{a very strong, very short-range attraction between nucleons.}
  \kw{Ionising radiation}{radiation that knocks electrons out of atoms, making ions.}
\end{keywords}

\begin{itemize}
  \item In $\chFourNuc{A}{Z}{X}$, $A$ = \term{mass number} (protons + neutrons) and $Z$ =
        \term{atomic number} (protons). Neutrons: $N = A - Z$. \term{Isotopes} have the same
        $Z$ but a different $N$.
  \item Protons repel each other (about \qty{230}{N} at \qty{1}{fm} $= \qty{1e-15}{m}$). The strong force
        attracts every nucleon more strongly, but only out to about \qty{3}{fm}. So in a big
        nucleus the repulsion grows faster. Extra neutrons add attraction but no repulsion.
  \item Stable nuclei lie in a thin \term{band of stability} ($N$ against $Z$): $N \approx Z$
        for light nuclei, extra neutrons for heavy ones. None with $Z > 82$ is stable.
  \item \term{Radioactive decay} is \term{spontaneous}: the \term{parent} nucleus changes by
        itself into the \term{daughter}, closer to the band. Alpha and beta decay change $Z$:
        a \term{transmutation} (one element turns into another).
  \item The beta particle is \emph{created} in the decay ($\beta^-$: a down quark becomes an
        up quark). A \term{positron} is the electron's antiparticle (charge $+e$). A
        \term{neutrino} is neutral, with almost no mass. A \term{gamma ray} is a very
        high-energy photon from an \term{excited} nucleus ($^{*}$).
\end{itemize}
\begin{align*}
  \alpha:\ & \chFourNuc{A}{Z}{X} \to \chFourNuc{A-4}{Z-2}{Y} + \chFourAlpha &
  \beta^-:\ & \chFourNuc{A}{Z}{X} \to \chFourNuc{A}{Z+1}{Y} + \chFourBeta + \chFourAnu\\
  \gamma:\ & \chFourNuc{A}{Z}{X}^{*} \to \chFourNuc{A}{Z}{X} + \chFourGam &
  \beta^+:\ & \chFourNuc{A}{Z}{X} \to \chFourNuc{A}{Z-1}{Y} + \chFourPos + \chFourNu
\end{align*}
\textbf{Balancing} (not on the formula sheet): nucleons and charge are conserved, so the
totals of $A$ and of $Z$ are the same on both sides. \Cref{ch4:tab:decays} compares the decays.

\begin{table}
  \caption{Alpha, beta and gamma decay and their radiation (u = atomic mass
    unit).}\label{ch4:tab:decays}
  \centering\small
  \begin{tabular}{@{}L{20mm}L{28mm}L{28mm}L{28mm}L{26mm}@{}}
    \toprule
    & \colhead{Alpha ($\alpha$)} & \colhead{Beta-minus ($\beta^-$)}
      & \colhead{Beta-plus ($\beta^+$)} & \colhead{Gamma ($\gamma$)}\\
    \midrule
    What happens & 2 protons and 2 neutrons leave together & a neutron turns into a proton
      & a proton turns into a neutron & an excited nucleus loses energy\\
    Emitted & helium-4 nucleus $\chFourAlpha$ & electron $\chFourBeta$ and antineutrino $\chFourAnu$
      & positron $\chFourPos$ and neutrino $\chFourNu$ & photon $\chFourGam$\\
    Change & $A - 4$, $Z - 2$ & $A$ same, $Z + 1$ & $A$ same, $Z - 1$ & none\\
    Which nuclei & mostly very heavy ($Z > 82$) & too many neutrons (above the band)
      & too few neutrons (below the band) & just after $\alpha$ or $\beta$ decay\\
    \midrule
    Charge & $+2e$ & $-e$ & $+e$ & 0\\
    Mass & about \qty{4}{\amu} & \qty{0.00055}{\amu} & \qty{0.00055}{\amu} & 0\\
    Speed & about 5\,\% of $c$ & up to 99\,\% of $c$ & up to 99\,\% of $c$ & $c$\\
    Ionising & strong & medium & medium & weak\\
    Stopped by & paper, skin or a few cm of air & a few mm of aluminium
      & a few mm of aluminium & only reduced by thick lead\\
    In a field & small bend & large bend, opposite way to $\alpha$
      & large bend, same way as $\alpha$ & no bend\\
    \bottomrule
  \end{tabular}
\end{table}

\begin{itemize}
  \item Alpha (heavy, slow, charge $+2e$) ionises many atoms in a short path and loses its
        energy fast. So the most ionising radiation is the least \term{penetrating}.
  \item In a field $\alpha$ and $\beta^-$ bend opposite ways (opposite charges). Beta bends
        more (much less momentum); gamma goes straight (\cref{ch4:fig:radiation}).
\end{itemize}

\begin{figure}
  \centering
  \begin{tikzpicture}[x=1cm, y=0.8cm]
    % (a) penetration
    \filldraw[object] (-0.6,-0.45) rectangle (0,1.65);
    \node[diagramlabel, rotate=90] at (-0.3,0.6) {source};
    \fill[hscmuted!45] (1.4,-0.5) rectangle (1.46,1.75);
    \fill[hscmuted!70] (3.0,-0.5) rectangle (3.15,1.75);
    \fill[hscink!55] (4.9,-0.5) rectangle (5.4,1.75);
    \draw[->, line width=1.1pt, hscforce] (0.1,1.3) -- (1.35,1.3);
    \node[diagramlabel, text=hscforce, above] at (0.55,1.3) {$\alpha$};
    \draw[->, line width=1.1pt, hscvel] (0.1,0.6) -- (2.95,0.6);
    \node[diagramlabel, text=hscvel, above] at (0.55,0.6) {$\beta$};
    \draw[line width=1.1pt, hscfield] (0.1,-0.1) -- (4.9,-0.1);
    \draw[->, line width=0.7pt, hscfield, dashed] (5.4,-0.1) -- (6.0,-0.1);
    \node[diagramlabel, text=hscfield, above] at (0.55,-0.1) {$\gamma$};
    \node[diagramlabel, align=center, anchor=north] at (1.43,-0.55) {paper\\\strut};
    \node[diagramlabel, align=center, anchor=north] at (3.08,-0.55) {aluminium\\(a few mm)};
    \node[diagramlabel, align=center, anchor=north] at (5.15,-0.55) {lead\\(several cm)};
    \node[diagramlabel, font=\small\bfseries] at (2.8,2.1) {(a) Penetration};
    % (b) deflection in a magnetic field into the page
    \begin{scope}[xshift=7.7cm, yshift=0.4cm]
      \foreach \p in {(0.6,1.2), (1.4,1.2), (2.2,1.2), (3.0,1.25),
          (0.6,0.5), (1.4,0.5), (3.0,0.45), (3.8,0.45),
          (0.6,-0.55), (1.4,-0.55), (2.2,-0.55), (3.0,-0.55), (3.8,-0.55)}
        \node[text=hscfield!70, font=\footnotesize] at \p {$\times$};
      \filldraw[object] (-0.7,-0.25) rectangle (-0.15,0.25);
      \draw[line width=1.1pt, hscforce, ->] (0,0) arc[start angle=-90, end angle=-58, radius=6];
      \node[diagramlabel, text=hscforce, right, fill=white, inner sep=1pt] at (3.25,0.95) {$\alpha$ (small bend)};
      \draw[line width=1.1pt, hscfield, ->] (0,0) -- (4.6,0);
      \node[diagramlabel, text=hscfield, right, fill=white, inner sep=1pt] at (4.6,0) {$\gamma$};
      \draw[line width=1.1pt, hscvel, ->] (0,0) arc[start angle=90, end angle=0, radius=1.1];
      \node[diagramlabel, text=hscvel, right, fill=white, inner sep=1pt] at (1.2,-1.1) {$\beta^-$ (large bend)};
      \node[diagramlabel, text=hscfield, anchor=north west] at (-0.7,-1.5)
        {$\vec{B}$ into the page ($\times$)};
      \node[diagramlabel, font=\small\bfseries] at (2.0,1.75) {(b) Deflection};
    \end{scope}
  \end{tikzpicture}
  \caption{(a) How far each radiation gets. (b) Paths in a magnetic field (not to scale).}
  \label{ch4:fig:radiation}
\end{figure}

%% ==================================================================
\section{Half-life and activity}\label{ch4:sec:halflife}
\covers{M8.IQ4.02}

\begin{bigidea}
One nucleus decays at a random moment. But in a huge sample, the amount halves in equal
time steps called half-lives.
\end{bigidea}

\begin{keywords}
  \kw{Half-life ($t_{\frac{1}{2}}$)}{the time for half of the radioactive nuclei in a sample to decay.}
  \kw{Decay constant ($\lambda$)}{the chance per second that one nucleus decays. Unit: \si{s^{-1}}.}
  \kw{Activity ($A$)}{the number of decays each second. Unit: becquerel (Bq).}
\end{keywords}

\begin{sidebyside}[0.40]
  \centering
  \begin{tikzpicture}
    \begin{axis}[hsc, width=0.95\linewidth, height=0.84\linewidth,
        xmin=0, xmax=4.6, ymin=0, ymax=1.1,
        xtick={0,1,2,3,4}, xticklabels={0, $t_{\frac{1}{2}}$, $2t_{\frac{1}{2}}$,
          $3t_{\frac{1}{2}}$, $4t_{\frac{1}{2}}$},
        ytick={0,0.25,0.5,1}, yticklabels={0, $N_0/4$, $N_0/2$, $N_0$},
        xlabel={Time $t$}, ylabel={Undecayed nuclei $N$}]
      \addplot[hscseriesA, line width=1.2pt, domain=0:4.6, samples=120, no marks] {2^(-x)};
      \draw[guide] (axis cs:1,0) -- (axis cs:1,0.5) -- (axis cs:0,0.5);
      \draw[guide] (axis cs:2,0) -- (axis cs:2,0.25) -- (axis cs:0,0.25);
      \draw[guide] (axis cs:3,0) -- (axis cs:3,0.125) -- (axis cs:0,0.125);
    \end{axis}
  \end{tikzpicture}
  \captionof{figure}{Exponential decay: the number of undecayed nuclei halves every
    half-life.}\label{ch4:fig:decay}
\nextside
  \begin{itemize}
    \item Decay is \term{random}: you cannot predict when one nucleus will decay.
    \item Each nucleus has the same chance each second. Age, temperature, chemistry and
          sample size do not change it.
    \item So the amount of the \emph{parent} halves every half-life (\cref{ch4:fig:decay}).
          After $n$ half-lives: $N_\mathrm{t} = N_0\left(\tfrac{1}{2}\right)^{n}$, with
          $n = t/t_{\frac{1}{2}}$.
    \item A \term{Geiger counter} measures a \term{count rate}. Minus the \term{background}
          (radiation always around us), it is proportional to the activity.
  \end{itemize}
\end{sidebyside}

\begin{formulacard}{The decay law and the decay constant}
  \cardformula{N_\mathrm{t} = N_0 e^{-\lambda t} \qquad\qquad
    \lambda = \frac{\ln 2}{t_{\frac{1}{2}}}}
  \begin{symbols}
    \sym{$N_\mathrm{t}$, $N_0$}{undecayed (parent) nuclei at time $t$; at the start}{no unit}
    \sym{$\lambda$}{decay constant}{\si{s^{-1}} (or \si{h^{-1}}, \si{\yr^{-1}})}
    \sym{$t$, $t_{\frac{1}{2}}$}{time since the start; half-life}{\si{s} (or h, yr)}
  \end{symbols}
  \whentouse{one isotope, many nuclei. Time: $t = \frac{1}{\lambda}\ln\frac{N_0}{N_\mathrm{t}}$.
    Also for mass or activity: $A_\mathrm{t} = A_0 e^{-\lambda t}$.}
  \trap{$\lambda$ and $t$ in one time unit; $\lambda$ is not a wavelength; use \texttt{ln}, not \texttt{log}.}
\end{formulacard}

\textbf{Activity} (not on the formula sheet): $A = \lambda N$. Each of the $N$ undecayed
nuclei has a chance $\lambda$ of decaying each second. For Bq, use $\lambda$ in \si{s^{-1}}.

%% ==================================================================
\section{Mass defect, binding energy and energy released}\label{ch4:sec:binding}
\covers{M8.IQ4.04, M8.IQ4.06}

\begin{bigidea}
A nucleus has less mass than its separate protons and neutrons. In any nuclear reaction
where the total mass goes down, the lost mass appears as energy ($E = mc^2$).
\end{bigidea}

\begin{keywords}
  \kw{Mass defect ($\Delta m$)}{the mass ``missing'' when nucleons join to make a nucleus.}
  \kw{Binding energy ($E_\mathrm{B}$)}{the energy needed to pull a nucleus apart into separate nucleons.}
\end{keywords}

\begin{itemize}
  \item When free nucleons join, energy is given out, so mass is lost. A nucleus is
        \emph{lighter} than its parts (helium-4: $\Delta m = \qty{0.030377}{\amu}$, or
        \qty{28.3}{MeV}).
  \item Binding energy must be \emph{supplied} to split a nucleus; it is not stored energy.
\end{itemize}

\begin{formulacard}[combine={$E = mc^2$ and $\mathrm{u} = 931.5\ \mathrm{MeV}/c^2$}]{Mass defect, binding energy and energy released}
  \cardformula{E = \Delta m\,c^2 \qquad E\,(\mathrm{MeV}) = \Delta m\,(\mathrm{u}) \times 931.5}
  \begin{symbols}
    \sym{$E$}{energy released (or the binding energy $E_\mathrm{B}$)}{\si{J} (or MeV)}
    \sym{$\Delta m$}{mass before $-$ mass after. For a nucleus (not on the sheet):
      $(Zm_\mathrm{p} + Nm_\mathrm{n}) - m_\mathrm{nucleus}$}{\si{kg} (or u)}
    \sym{$m_\mathrm{p}$, $m_\mathrm{n}$}{\qty{1.007276}{\amu}, \qty{1.008665}{\amu} (the sheet gives kg only)}{\si{kg} (or u)}
  \end{symbols}
  \whentouse{any nuclear reaction, and binding energy. Given binding energies (fission,
    fusion): $E = \sum E_\mathrm{B,after} - \sum E_\mathrm{B,before}$.}
  \derive{$E = mc^2$ with $\qty{1}{\amu} = \qty{931.5}{MeV}/c^2$ (\qty{1}{\amu} = \qty{1.661e-27}{kg}).}
  \trap{one route only: u $\times$ 931.5 gives MeV; kg $\times\, c^2$ gives J. Keep every
    decimal place. Atomic masses include electrons (\qty{0.000549}{\amu} each). For a mass
    defect with $m_\mathrm{p}$, subtract $Z \times \qty{0.000549}{\amu}$ from the atomic mass
    first. In a reaction (alpha, $\beta^-$, fission, fusion) the electrons cancel, so use
    atomic masses as they are; in $\beta^+$ decay also subtract $2m_\mathrm{e}$.}
\end{formulacard}

\textbf{Energy released in decays and transmutations.}
\begin{itemize}
  \item Mass-energy, $A$ and charge are conserved. Less mass after: the lost mass becomes
        kinetic energy and gamma rays. More mass after: energy must be supplied.
  \item \textbf{Alpha decay:} momentum stays zero, so the daughter \term{recoils}. The light
        alpha particle gets most of the energy (about \qty{98}{\percent}), so alpha energies
        are fixed.
  \item \textbf{Beta decay:} each decay releases the same energy, but beta energies are
        spread. Pauli (1930): a neutrino carries off the rest (first detected in 1956).
  \item \term{Artificial transmutation}: a change forced by firing a particle at a nucleus.
        Rutherford (1919, the first): $\chFourNuc{14}{7}{N} + \chFourAlpha \to
        \chFourNuc{17}{8}{O} + \chFourNuc{1}{1}{H}$. Cockcroft and Walton (1932):
        $\chFourNuc{7}{3}{Li} + \chFourNuc{1}{1}{H} \to 2\,\chFourAlpha$, about \qty{17}{MeV},
        as the mass loss predicted (the first direct test of $E = mc^2$).
\end{itemize}

\begin{sidebyside}[0.73]
  \centering
  \includegraphics{ch4_becurve.pdf}
  \captionof{figure}{Binding energy per nucleon against mass number. Fission and fusion
    both move nuclei up the curve.}\label{ch4:fig:becurve}
\nextside
  \small
  \begin{itemize}
    \item \term{Binding energy per nucleon} $= E_\mathrm{B}/A$ (\cref{ch4:fig:becurve}).
          Bigger means more tightly bound.
    \item Hydrogen-2: only \qty{1.1}{MeV}.
    \item Helium-4: unusually high (\qty{7.1}{MeV}).
    \item Peak near iron-56 (\qty{8.8}{MeV}): the most tightly bound.
    \item Uranium: about \qty{7.6}{MeV}.
    \item \textbf{Key rule:} moving up the curve, towards iron, means less mass, so energy
          is released.
  \end{itemize}
\end{sidebyside}

%% ==================================================================
\section{Fission and fusion}\label{ch4:sec:fission}
\covers{M8.IQ4.03, M8.IQ4.05, M8.IQ4.04, M8.IQ4.06}

\begin{bigidea}
Splitting a heavy nucleus (\term{fission}) and joining light nuclei (\term{fusion}) both
make nucleons more tightly bound. The products have less mass, so energy is released.
\end{bigidea}

\begin{keywords}
  \kw{Chain reaction}{neutrons from one fission cause more fissions.}
  \kw{Critical mass}{the smallest mass of fuel that can keep a chain reaction going.}
\end{keywords}

% keep this lead-in line on the same page as its list (layout check: it was alone at a
% page bottom)
\Needspace{4\baselineskip}
\textbf{How U-235 splits} (Hahn and Strassmann, 1938):
\begin{enumerate}
  \item It absorbs a slow (\term{thermal}) neutron. (A neutron has no charge, so it is not
        repelled.) It becomes an excited U-236 and stretches into a peanut shape.
  \item The short-range strong force cannot hold the far ends. Repulsion pushes them apart:
        two fragments (their kinetic energy becomes heat) and 2 or 3 fast neutrons.
\end{enumerate}
\[
  \chFourNuc{235}{92}{U} + \chFourNeut \to \chFourNuc{236}{92}{U}^{*} \to
  \chFourNuc{140}{54}{Xe} + \chFourNuc{94}{38}{Sr} + 2\,\chFourNeut \qquad
  \text{(one possible pair)}
\]
\textbf{Energy:} U-235 has \qty{7.6}{MeV} per nucleon; the fragments have 8.3 to
\qty{8.6}{MeV}. More tightly bound means less mass: about \qty{200}{MeV} is released per
fission. The fragments are radioactive (nuclear waste).

\textbf{Chain reactions.} The 2 or 3 new neutrons can cause more fissions. What matters is
how many neutrons from each fission cause another fission, on average (\cref{ch4:fig:chain}):
\begin{itemize}
  \item \emph{more than one}: fissions multiply (1, 3, 9, 27, \ldots). This
        \term{uncontrolled chain reaction} is a nuclear explosion;
  \item \emph{exactly one}: a steady rate. This \term{controlled chain reaction} runs a
        reactor (first done by Fermi's team, 1942);
  \item \emph{less than one}: the chain dies out.
\end{itemize}

\begin{figure}
  \centering
  \begin{tikzpicture}[x=1cm, y=0.52cm, nuc/.style={circle, draw=hscink, fill=hscacc!25,
      line width=0.6pt, inner sep=0pt}, neu/.style={->, hscmuted, line width=0.7pt}]
    \node[nuc, minimum size=0.34cm] (a0) at (0,0) {};
    \foreach \i/\y in {1/1.5, 2/0, 3/-1.5} {
      \node[nuc, minimum size=0.28cm] (a\i) at (2,\y) {};
      \draw[neu] (a0) -- (a\i);
      \foreach \j/\dy in {1/0.5, 2/0, 3/-0.5} {
        \node[nuc, minimum size=0.2cm] (b\i\j) at (3.9,\y+\dy) {};
        \draw[neu] (a\i) -- (b\i\j);
      }
    }
    \node[diagramlabel] at (0,-2.9) {1 fission};
    \node[diagramlabel] at (2,-2.9) {3 fissions};
    \node[diagramlabel] at (3.9,-2.9) {9 fissions};
    \node[diagramlabel, font=\small\bfseries] at (2,2.9) {(a) Uncontrolled: more than one};
    \begin{scope}[xshift=6.3cm]
      \fill[hscink!65] (-0.2,1.3) rectangle (6.3,1.7);
      \node[diagramlabel, above] at (3.05,1.7) {control rod absorbs neutrons};
      \foreach \k in {0,1,2,3} \node[nuc, minimum size=0.34cm] (c\k) at (\k*1.9,0) {};
      \foreach \k/\n in {0/1,1/2,2/3} {
        \draw[neu] (c\k) -- (c\n);
        \draw[neu] (c\k) -- ++(0.55,1.25);
        \draw[neu] (c\k) -- ++(0.55,-1.2);
      }
      \draw[neu] (c3) -- ++(0.55,1.25);
      \draw[neu] (c3) -- ++(0.55,-1.2);
      \node[diagramlabel, text=hscmuted] at (3.05,-1.8) {other neutrons escape};
      \node[diagramlabel] at (3.05,-2.9) {1 fission $\to$ 1 fission $\to$ 1 fission $\to \cdots$};
      \node[diagramlabel, font=\small\bfseries] at (3.05,2.9) {(b) Controlled: exactly one};
    \end{scope}
  \end{tikzpicture}
  \caption{Chain reactions. Each arrow is a neutron; each circle is a U-235 nucleus that
    splits.}\label{ch4:fig:chain}
\end{figure}

\begin{itemize}
  \item A small lump loses many neutrons through its surface. Below the critical mass
        (\term{subcritical}) the chain dies out; above it (\term{supercritical}) it grows.
  \item Natural uranium is only \qty{0.7}{\percent} U-235. Reactor fuel is \term{enriched} to
        3 to \qty{5}{\percent} U-235; bombs need about \qty{90}{\percent}.
  \item \textbf{Reactor:} the \term{moderator} (water or graphite) \emph{slows} neutrons, so
        they split U-235 more easily. \term{Control rods} (boron or cadmium) \emph{absorb}
        neutrons, so exactly one per fission causes a new fission. The \term{coolant} carries
        the heat away to make steam.
  \item \textbf{Bomb:} highly enriched fuel is pushed together into a supercritical mass.
        Nothing absorbs the extra neutrons, so the energy comes out in a fraction of a second.
\end{itemize}

\textbf{Nuclear fusion.}
\begin{itemize}
  \item \textbf{Why energy is released:} helium-4 (\qty{7.1}{MeV} per nucleon) is far more
        tightly bound than hydrogen-2 (\qty{1.1}{MeV}), so the products have less mass.
  \item \textbf{Why it is hard:} both nuclei are positive and repel. They must get within
        about \qty{3}{fm} for the strong force to act, so the gas must be very hot and dense.
        The Sun's core is about \qty{1.5e7}{K}; reactors on Earth need over \qty{1e8}{K}.
  \item The Sun: $4\,\chFourNuc{1}{1}{H} \to \chFourAlpha + 2\,\chFourPos + 2\,\chFourNu$,
        \qty{26.7}{MeV} in total (with positron annihilation).
  \item Future reactors: $\chFourNuc{2}{1}{H} + \chFourNuc{3}{1}{H} \to \chFourAlpha +
        \chFourNeut$ (\qty{17.6}{MeV}), using \term{deuterium} (hydrogen-2: one proton, one
        neutron) and \term{tritium} (hydrogen-3: one proton, two neutrons).
  \item Fusing nuclei beyond iron would take in energy.
\end{itemize}

%% ==================================================================
\section{Key experiments}\label{ch4:sec:experiments}
\covers{M8.IQ4.01, M8.IQ4.02}

\textbf{Dice model of half-life:} throw 200 dice and remove every six; repeat with the rest.
The number left falls like \cref{ch4:fig:decay}: half are left after about 3.8 throws,
whatever the starting number.

\begin{experiment}{Penetrating power of alpha, beta and gamma radiation}
\textbf{Aim:} to compare how well $\alpha$, $\beta$ and $\gamma$ pass through materials.
\textbf{What was done:} the background count was measured first. Each sealed source (Am-241:
$\alpha$; Sr-90: $\beta$; Co-60: $\gamma$) was put a fixed distance from a Geiger--M\"uller
tube. Counts were taken with no absorber, paper, \qty{5}{mm} of aluminium and \qty{10}{mm} of
lead. \textbf{Result:} paper stops alpha; \qty{5}{mm} of aluminium stops beta; \qty{10}{mm} of
lead only halves the gamma count. \textbf{Conclusion and significance:} penetrating power
$\gamma > \beta > \alpha$; ionising power is the reverse. This guides shielding and uses: alpha
in smoke alarms, gamma in cancer treatment.
\end{experiment}

\begin{summarybox}{Remember}
\begin{itemize}
  \item Short-range strong force against proton repulsion: big nuclei, or the wrong mix of
        $N$ and $Z$, are unstable.
  \item $\alpha$: $A - 4$, $Z - 2$. $\beta^-$: n $\to$ p. $\beta^+$: p $\to$ n. $\gamma$:
        energy only. Ionising $\alpha > \beta > \gamma$; penetrating the reverse.
  \item $N_\mathrm{t} = N_0 e^{-\lambda t}$, $\lambda = \ln 2/t_{\frac{1}{2}}$,
        $A = \lambda N$ (Bq: $\lambda$ in \si{s^{-1}}).
  \item $E = \Delta m\,c^2$, or $\Delta m$ (u) $\times$ 931.5 MeV. Fission and fusion move up
        the curve towards iron-56: less mass, energy out.
  \item Alpha energies are fixed; beta energies are spread (the neutrino). Reactor: exactly one
        neutron per fission causes a fission; moderator slows, control rods absorb.
\end{itemize}
\end{summarybox}
